\scanpage{5}
\mspage{L 16}

suis dans \handwritten{\inserted{les pages de}} Récoltes et Semailles. \inserted{Et} \deleted{\unsure{A}}aucun des trois n'est a\unsure{chevé|bhevé} avec le point final\handwritten{,} à la page douze cents et quelques. Les échos, déjà, que \unsure{ma} recueillir mon témoignage (et jusques y compris l'écho par le silence \ldots), \deleted{\unsure{fera}}f\handwritten{eront} partie de la ``suite'' du voyage. \deleted{\unsure{Mais à vrai dire,|Mais, à vrai dire,}}\inserted{Quant à \unsure{un|son} ``terme'',} ce voyage \inserted{sûrement} est de ceux qui ne sont jamais menés à terme - pas même, si ça se trouve, au jour de notre mort.

Et me \deleted{\unsure{re}}voilà \deleted{revenu} enfin \handwritten{\inserted{revenu}} au point de départ: te dire d'avance, si faire se peut, ``de quoi il est question'' dans Récoltes et Semailles. Mais il est vrai aussi, que sans l'avoir cherché, les pages précédentes te l'ont déjà dit peu ou prou. Il sera plus intéressant, peut-être, de continuer \deleted{plutôt} sur ma lancée, \handwritten{\inserted{et}} de \unsure{raconter|raonter}, plutôt que d'annoncer.

Juin 1985\\

\deleted{J'ai commencé à} Les page\unsure{s} qui précèdent de cette lettre (ou ``pré-lettre\unsure{'|''} ont été écrites \deleted{à un moment creux} le mois dernier, \unsure{en|an} un court ``moment creux''. Entretemps, j'ai fini de mettre la dernière main à la quatrième partie de Récoltes et Semailles (``Les Quatre Opérations'') - il ne reste plus que cette lettre à terminer, pour que tout soit prêt enfin pour la frappe et pour la duplication. Depuis \inserted{bientôt} un an et demi que je suis``sur le point de terminer'' ces fameuses notes,\deleted{\unsure{eux|aux} et \unsure{sa|son} introduction,} je n'y croyais plus \ldots

En me mettant, au mois de février l'an dernier, à cette ``introduction'' de nature un peu inhabituelle à un ouvrage mathématique, il y avait (je crois) trois genres de choses surtout sur lesquelles j'avais envie de m'exprimer. D'une part, \deleted{\unsure{comme}} m'expliquer sur mes intentions en revenant à une activité mathématique, et sur l'esprit dans lequel j'avais écrit ce premier volume de ``A la Poursuite des Champs'', et dans lequel je comptais poursuivre un vaste voyage de prospection et de découverte mathématiques, dans les ``Réflexions''.\deleted{(} Il ne s'agirait plus, désormais, de présenter des fondements méticuleux et à quatre épingles de théories en gésine, mais de ``carnets de bord'' plutôt, où \unsure{le|un} travail \deleted{de prospection} se poursuivrait au jour le jour, tel qu'il se fait vraiment, depuis les premières intuitions ou idées de départ, floues souvent \deleted{voire \unsure{à}} et réticentes à se laisser saisir dans les mailles du langage,\deleted{(}alors que c'est le langage adéquat à les saisir qui souvent fait justement défaut et qu'il s'agit de faire se condenser d'un apparent néant de brumes \deleted{parfois} impalpables \ldots{} Depuis ce\deleted{\unsure{s}} stade\deleted{\unsure{s}} donc de ce qui n'est encore que pressenti avant d'être

% archive sheet number: 4
